% Section 2: Background and Taxonomy
\section{Background and Taxonomy}

\subsection{REST vs. Prompt-Based APIs}

\gls{rest} defines architectural constraints for distributed systems: statelessness (requests contain all necessary information) and uniform interface (standard \gls{http} verbs with typed schemas) \cite{FIELDING_2000}. \gls{rest} \gls{api}s excel at \gls{crud} operations with deterministic, validated responses \cite{RICHARDSON_2007}.

Prompt-based \gls{api}s violate these principles fundamentally. Natural language lacks fixed schemas; the same intent can be expressed infinitely. Statelessness conflicts with conversational context. Yet this flexibility, enabled by transformer architectures \cite{VASWANI_2017}, provides unprecedented interface simplicity \cite{BROWN_2020}. Table \ref{tab:api_comparison} contrasts key properties.

\begin{table}[h]
\centering
\caption{Traditional REST vs. Prompt-Based APIs}
\label{tab:api_comparison}
\small
\begin{tabular}{|p{2.2cm}|p{2.5cm}|p{2.5cm}|}
\hline
\textbf{Aspect} & \textbf{Traditional REST} & \textbf{Prompt-Based} \\
\hline
Input & Structured JSON & Natural language \\
\hline
Validation & Schema-enforced & Semantic \\
\hline
Determinism & Yes & No \\
\hline
State & Stateless & Often contextual \\
\hline
Testing & Unit tests & LLM-as-judge \\
\hline
Latency & $<10ms$ & $100-2000ms$ \\
\hline
Cost & Compute/storage & Pay-per-token \\
\hline
\end{tabular}
\end{table}

\subsection{Taxonomy of Prompt-Based Services}

\begin{enumerate}
    \item \textbf{Direct Prompt:} Single-turn, stateless interactions for classification or extraction tasks.
    \item \textbf{Conversational:} Multi-turn interactions requiring context management, e.g., customer support.
    \item \textbf{Agentic:} LLMs with tool access, necessitating secure function calling and sandboxing \cite{AWS_BEDROCK, SCHLUNTZ_2024}.
    \item \textbf{Hybrid:} Combining REST for transactional operations with prompts for flexible tasks, e.g., using REST for order processing and prompts for customer queries \cite{LERCHER_2024, CHEN_2025}.
\end{enumerate}

\subsection{AWS Bedrock Platform}

Bedrock provides managed foundation model access via unified \gls{api} \cite{AWS_BEDROCK}. \textbf{No free tier; all inference is pay-per-use.} Nova Micro costs \$0.000035 per 1K input tokens, the lowest on Bedrock \cite{AWS_NOVA}. Bedrock's Converse \gls{api} normalizes provider-specific formatting, enabling model swapping without client changes, which is critical for flexible service design \cite{AWS_BEDROCK}.
